\subsection{Representación por bloques del vector de estado}

La decisión estructural principal consiste en particionar el vector de estado en bloques independientes de tamaño fijo. Cada bloque agrupa un número
determinado de estados base consecutivos y almacena sus partes real e imaginaria como un arreglo contiguo de números en punto flotante. Esta
organización convierte el problema global de compresión del estado completo en una colección de problemas locales sobre subarreglos más pequeños.

La ventaja de esta representación es doble. Por una parte, cada bloque puede comprimirse y descomprimirse sin depender del contenido de los demás,
lo que facilita actualizaciones localizadas. Por otra, el tamaño del bloque se convierte en un parámetro de diseño que permite balancear dos efectos
opuestos: bloques más grandes tienden a favorecer el ratio de compresión, mientras que bloques más pequeños reducen el costo de acceso cuando solo
una fracción del vector es requerida.

\begin{figure}[H]
    \centering
    \fbox{\parbox{0.82\linewidth}{\centering
        Espacio reservado para una figura del vector de estado particionado en bloques, indicando el mapeo entre índices de estados base y bloques
        comprimidos.}}
    \caption{Particionamiento abstracto del vector de estado en bloques independientes}
    \label{fig:simulation-design-blocks}
\end{figure}
